% !TeX root = ../凸优化笔记.tex
% 由主文档通过 \input 载入；请编译主文档。
\setcounter{chapter}{5} % 保留原书章号。
\chapter{逼近、拟合与正则化}
\label{ch:co-fitting}

拟合问题需要回答两件事：什么样的误差不可接受，以及什么样的解更符合任务需要。前者决定损失函数，后者表现为约束或正则项。本章固定设计矩阵 $A\in\mathbb R^{r\times n}$ 与观测 $b\in\mathbb R^r$，以 $e(x)=Ax-b$ 记残差。

\section{范数逼近与误差度量}

最小二乘最小化 $\sum_i e_i^2$，对大残差施加较重惩罚；最小绝对偏差最小化 $\sum_i|e_i|$，对残差大小的惩罚为线性；极小极大逼近最小化 $\max_i|e_i|$，控制最坏单点误差。后两者的 LP 表示已在第\ref{ch:co-models}章给出。

最小化 $\|e\|_2$ 与 $\|e\|_2^2$ 有相同最优解，因为平方在非负实数上严格递增。不同范数一般没有这种对应。权重模型 $\sum_iw_i e_i^2$ 在 $w_i\geq0$ 时凸；权重可以表达量纲归一化或已知的观测可靠性，但应与建模假设相符。

平方损失对应的单点导数是 $2e_i$，残差越大，对更新方向的影响越大；绝对值损失的次梯度大小至多为一，因此单个大残差的影响受到限制。极小极大准则更关心目前最大的那几个误差，允许其他较小误差适当增大。选择损失之前，应先说清楚任务究竟关心平均表现、抗异常能力，还是最坏情况，不能仅凭某个模型容易计算来决定。

\begin{example}[均值、中位数与中程数]
用常数 $x$ 拟合观测 $0,0,9$。平方损失的唯一解为均值 $3$；绝对偏差损失为 $2|x|+|x-9|$，在 $x=0$ 唯一最小；最大绝对误差要求同时接近最小值 $0$ 和最大值 $9$，故最优解为 $9/2$，误差为 $9/2$。同一数据下，不同目标表达了不同偏好。
\end{example}

均值、中位数和中程数的差别也可以一般地推出。对有序数据 $b_{(1)}\leq\cdots\leq b_{(N)}$，绝对偏差和在不等于观测值的位置的导数为“左侧观测个数减右侧观测个数”。导数从负变正的位置就是中位数；若 $N$ 为偶数，两个中间观测之间的整个区间都最优。对于最大绝对偏差，只有两端 $b_{(1)},b_{(N)}$ 决定最坏距离，二者的中点给出最小半径 $(b_{(N)}-b_{(1)})/2$。这些结论说明非唯一性有时来自损失的平坦段，并不是算法没有收敛。

\section{最小范数解与带约束逼近}

欠定方程 $Ax=b$ 可能有很多解。最小范数模型 $\min\|x\|_2^2/2$、$Ax=b$ 选出唯一最小欧氏范数解；当 $A$ 满行秩时，第\ref{ch:co-duality}章已得
\[
 x^\star=A^{\mathsf T}(AA^{\mathsf T})^{-1}b.
\]
任意其他解为 $x^\star+d$，$Ad=0$；因 $x^\star\in\operatorname{range}A^{\mathsf T}$，有 ${x^\star}^{\mathsf T}d=0$，故 $\|x^\star+d\|_2^2=\|x^\star\|_2^2+\|d\|_2^2$，再次证明唯一性。一般可行的秩亏系统可用伪逆 $A^\dagger b$ 表示，见附录\ref{app:co-numerical}。

实际拟合常要求 $x\geq0$、$\mathbf1^{\mathsf T}x=1$ 或 $l\leq x\leq u$。这些约束仿射，故与凸损失组合后仍凸。例如非负最小二乘的 KKT 条件为 $A^{\mathsf T}(Ax-b)-\lambda=0$、$x,\lambda\geq0$、$x_i\lambda_i=0$。将无约束解的负分量直接截断通常不是其解，因为矩阵 $A^{\mathsf T}A$ 会耦合各分量。

\begin{example}[为什么不能总是先拟合再截断]
令 $A=\left[\begin{smallmatrix}1&1\\0&1\end{smallmatrix}\right]$、$b=(0,1)^{\mathsf T}$，求非负最小二乘解。
\end{example}
\begin{solve}
无约束方程 $Ax=b$ 的解为 $(-1,1)$，逐分量截断给 $(0,1)$，其目标为 $1/2$。在非负域上，固定 $x_2\geq0$ 时，第一项 $(x_1+x_2)^2/2$ 在 $x_1=0$ 最小；于是问题化为最小化 $[x_2^2+(x_2-1)^2]/2$，解为 $x_2=1/2$，目标为 $1/4$。真正的最优解是 $(0,1/2)$。第一坐标碰到边界后，第二坐标也必须重新调整，不能保持原无约束值不动。
\end{solve}

\section{正则化与偏好编码}

正则化（regularization）在数据损失之外加入对解的偏好。给定 $\alpha>0$，两个典型模型为
\begin{equation}
 \min_x\frac12\|Ax-b\|_2^2+\frac\alpha2\|x\|_2^2,
 \qquad
 \min_x\frac12\|Ax-b\|_2^2+\alpha\|x\|_1.
 \label{eq:co-regularized-fitting}
\end{equation}
前者为岭回归（ridge regression），其唯一解为 $(A^{\mathsf T}A+\alpha I)^{-1}A^{\mathsf T}b$。后者为 Lasso，最优性条件为
\[
 0\in A^{\mathsf T}(Ax-b)+\alpha\partial\|x\|_1.
\]
因此当最优分量 $x_j=0$ 时，相应残差相关量可在 $[-\alpha,\alpha]$ 内变化；这解释了绝对值正则项产生零分量的机制，但不保证任意数据都能恢复真实稀疏结构。

岭回归的稳定作用可以用奇异值分解看清。若 $A=U\Sigma V^{\mathsf T}$ 的正奇异值为 $\sigma_1,\ldots,\sigma_r$，相应左右奇异向量为 $u_j,v_j$，则
\[
 x_\alpha=\sum_{j=1}^r
 \frac{\sigma_j}{\sigma_j^2+\alpha}(u_j^{\mathsf T}b)v_j.
\]
未正则化的最小范数解使用系数 $1/\sigma_j$，很小的奇异值可能显著放大数据扰动；岭回归把它换成 $\sigma_j/(\sigma_j^2+\alpha)$，从而压低这些方向的响应。代价是即使数据没有噪声，也通常不再精确拟合。正则化改善稳定性与引入偏差，是同一次模型选择的两个方面。

\begin{example}[软阈值]
当 $A=I$ 时，Lasso 可逐坐标求解 $\min_x(x-b)^2/2+\alpha|x|$。在正半轴驻点为 $b-\alpha$，负半轴驻点为 $b+\alpha$；零点最优当且仅当 $|b|\leq\alpha$。故
\[
 x^\star=\operatorname{sign}(b)\max\{|b|-\alpha,0\},
\]
其中 $\operatorname{sign}(0)=0$，称为软阈值（soft thresholding）。对 $b=(3,1/2)$、$\alpha=1$，解为 $(2,0)$。
\end{example}

若 $x_\alpha$ 最小化 $L(x)+\alpha R(x)$，则它也最小化约束问题 $\min L(x)$、$R(x)\leq R(x_\alpha)$，否则可找到更小的惩罚目标。反向对应依赖乘子存在性，且参数不必一一对应。由两个权重下的最优性不等式相加可知，随 $\alpha$ 增大，选取的最优解的正则项值不增，数据损失值不减。

具体地，设 $0<\alpha<\beta$ 且两问题均有有限目标的最优解，则
\[
 L(x_\alpha)+\alpha R(x_\alpha)\leq L(x_\beta)+\alpha R(x_\beta),
 \qquad
 L(x_\beta)+\beta R(x_\beta)\leq L(x_\alpha)+\beta R(x_\alpha).
\]
相加得到 $(\beta-\alpha)[R(x_\beta)-R(x_\alpha)]\leq0$，再代回第一式得 $L(x_\alpha)\leq L(x_\beta)$。这个结论只描述训练模型内部的权衡，并不说明未见数据上的误差一定怎样变化。实际选取 $\alpha$ 时，可用独立验证数据比较预测损失；用来选择参数的数据与最终评价数据应在概念上区分。

\section{鲁棒损失与不确定性集合}

Huber 损失在小误差处为平方、大误差处为线性。给定 $\delta>0$，定义
\[
 h_\delta(t)=
 \begin{cases}t^2/2,&|t|\leq\delta,\\
 \delta|t|-\delta^2/2,&|t|>\delta.
 \end{cases}
\]
其导数为截断到 $[-\delta,\delta]$ 的 $t$，连续且单调不减，故损失凸。还可写成 $h_\delta(t)=\min_v\{(t-v)^2/2+\delta|v|\}$，由软阈值代入验证。这使大残差的边际惩罚不再无限增长。

辅助变量 $v$ 可以解释为由异常因素承担的那部分残差，小残差留在平方项里，大残差则部分转交给绝对值项。于是 Huber 拟合等价于联合凸问题
\[
 \min_{x,v}\ \frac12\|Ax-b-v\|_2^2+\delta\|v\|_1.
\]
当残差绝对值不超过 $\delta$ 时，最优 $v_i=0$；超过时，$v_i$ 承担超出的部分。这个表示既推导了损失的分段公式，也给出了另一种可实现的求解方式。

鲁棒优化（robust optimization）则要求约束对指定不确定性集合内的所有参数成立。设约束系数为 $a=\bar a+Pu$，$\|u\|_2\leq\rho$，其中 $P$ 为已知矩阵、$\rho\geq0$。条件 $a^{\mathsf T}x\leq b$ 对所有这样的 $u$ 成立，当且仅当
\[
 \bar a^{\mathsf T}x+\rho\|P^{\mathsf T}x\|_2\leq b.
\]
因为 $\sup_{\|u\|_2\leq\rho}u^{\mathsf T}P^{\mathsf T}x=\rho\|P^{\mathsf T}x\|_2$，上界由柯西--施瓦茨不等式给出，平行方向取到。所得为 SOCP 约束；若不确定集改为 $\|u\|_\infty\leq\rho$，则对应项为 $\rho\|P^{\mathsf T}x\|_1$。

最坏情形目标 $\sup_{u\in U}F(x,u)$ 在每个 $F(\cdot,u)$ 凸时仍凸。保证的范围由集合 $U$ 决定；抗异常损失、平均风险和最坏情形约束表达的是不同假设。

\section{函数拟合、平滑与插值}

给定固定基函数 $\phi_1,\ldots,\phi_n$，令拟合函数为 $g(t)=\sum_jx_j\phi_j(t)$。在观测位置 $t_i$ 上有 $g(t_i)=(Ax)_i$，其中 $A_{ij}=\phi_j(t_i)$，所以损失关于系数 $x$ 的凸性与基函数关于 $t$ 是否凸没有直接关系。

对网格上的函数值 $x_1,\ldots,x_n$，一阶差分矩阵 $D$ 满足 $(Dx)_i=x_{i+1}-x_i$。平方平滑项 $\|Dx\|_2^2$ 惩罚快速变化；总变差（total variation）$\|Dx\|_1$ 倾向产生较少的跳变。固定网格 $t_1<\cdots<t_n$ 上，分段线性插值函数凸当且仅当相邻斜率满足
\[
 \frac{x_{i+1}-x_i}{t_{i+1}-t_i}
 \leq\frac{x_{i+2}-x_{i+1}}{t_{i+2}-t_{i+1}},\quad i=1,\ldots,n-2.
\]
这些是关于函数值的线性不等式。对于任意其他函数族，仅在采样点检验形状通常不能认证整个连续区间上的性质。

精确插值要求 $Ax=b$，可能在噪声下放大不稳定性；逼近允许残差并通过损失和正则项进行权衡。两者首先是不同的模型选择。

\begin{example}[由单调性约束修正一组观测]
在三个依次增大的位置观测到 $b=(3,1,2)$，认为真实数值应当不减，求 $\min_x\|x-b\|_2^2/2$，满足 $x_1\leq x_2\leq x_3$。
\end{example}
\begin{solve}
前两个观测的顺序冲突，把它们调整为共同均值 $2$，得到候选点 $x=(2,2,2)$。对约束 $x_1-x_2\leq0$、$x_2-x_3\leq0$ 配乘子 $\lambda_1,\lambda_2$，驻点式为
\[
 x-b+(\lambda_1,-\lambda_1+\lambda_2,-\lambda_2)^{\mathsf T}=0.
\]
取 $(\lambda_1,\lambda_2)=(1,0)$ 即满足全部 KKT 条件，因此该候选点最优；平方目标严格凸，解唯一。第二条约束虽然活跃，乘子却为零，这也具体展示了活跃性与正乘子的区别。
\end{solve}

\section{综合例题与学习检查}
\begin{enumerate}
 \item 对观测 $0,0,M$，分别求三种范数下的最佳常数，分析 $M\to+\infty$ 时的变化。答案为 $M/3,0,M/2$。
 \item 推导带非负约束的岭回归 KKT 条件。说明为何 $A=I$ 时可以逐坐标求解，而一般 $A$ 不行。
 \item 给定线性拟合的椭球系数不确定性，写出其 SOCP 鲁棒对应形式，并核对矩阵维度。
 \item 在等距网格上拟合一列数据，写出“平方误差加总变差”和“平方误差加凸性约束”两个模型。前者偏好分段常数，后者限制斜率单调，两种模型不可互换。
 \item 推导练习：对 $A=I$、$b=(3,1/2)$，写出岭回归与 Lasso 的解随 $\alpha>0$ 变化的公式，比较哪些坐标会严格变成零。答案分别为 $b/(1+\alpha)$ 与逐坐标软阈值。
\end{enumerate}

